\begin{proposition}[Cotangent direct images and Hodge numbers of the covers]\hypertarget{tag-01S8}{}\label{inv:B3.2}
\StatementRecord{inv:B3.2}{proved}{01S8}{\SeriesRef{inv:B3.1}, \SeriesRef{inv:B1.2}}

Write $T_n=\C_{p_0}^{\,n-1}$. Then
\[
\begin{array}{c|ccc}
 &R^0g_*&R^1g_*&R^2g_*\\\hline
\Omega^1_Y&\OO(-2)&\OO\oplus\OO(-1)^2\oplus T_n&\OO\oplus\OO(-1)\\
\Omega^2_Y&\OO(-2)\oplus\OO(-1)&\OO(-2)\oplus\OO(-1)^2&\OO\oplus T_n
\end{array}
\]
Writing $V$ for the locally free quotient, the splittings form an
affine space under $\operatorname{Hom}(V,T_n)$, by the exact
sequence in the proof.
The only nonzero Hodge numbers of $Y$ are
\[
h^{0,0}=h^{0,1}=h^{3,2}=h^{3,3}=1,\quad
h^{1,2}=h^{2,1}=1,\quad h^{1,1}=h^{2,2}=n+1.
\]
The original finite-fundamental-group threefold $X$ itself has the
same Hodge diamond as the simply connected $(3,4)$ member.
\end{proposition}


\begin{proof}
\begin{enumerate}[label=\arabic*.,leftmargin=*]
\item \ref{inv:B3.1} identifies the logarithmic cotangent sequence on cohomology,
and therefore its connecting maps and direct images, with those of
the original family. The normalized residue sheaf differs only at
the cusp: its degree-zero cohomology there is now $\C^n$.
The residue of a base logarithmic form is the diagonal vector
$(1,\ldots,1)$, since each cusp component has multiplicity one.
Its cokernel is $T_n$. The degree-one residues at the finite fibres
cancel the same two skyscrapers as in \ref{inv:B1.2}. The long exact sequence
therefore gives
\[
0\longrightarrow T_n\longrightarrow R^1g_*\Omega^1_Y
\longrightarrow\OO\oplus\OO(-1)^2\longrightarrow0
\]
and the remaining entries of the first row. The extension splits
because $\operatorname{Ext}^1(V,T_n)=H^1(V^*\otimes T_n)=0$
for a vector bundle $V$ on $B$.

\item For a sheaf $F$ locally free on the total space, relative duality
in its derived form gives
\[
Rg_*(F^*\otimes\omega_{Y/B})
\simeq R\mathcal Hom_B(Rg_*F,\OO_B)[-2].
\]
For $F=\Omega^1_Y$, the torsion in $R^1g_*F$ contributes
$\mathcal Ext^1_B(T_n,\OO_B)$ in degree two of the right side.
This is a length-$(n-1)$ skyscraper at $p_0$. After tensoring by
$\OO(-2)$ it still has that length. The locally free parts dualize
in complementary degrees, giving exactly the second row; the
extension by the locally free quotient splits for the same reason
as above. Coherent Leray computes the displayed Hodge numbers.
For $X$ with one-component cusp, \ref{inv:B1.2} already applies independently
of its fundamental group.
\end{enumerate}
\end{proof}